% Generated draft: requires mathematical and source review.
\documentclass[11pt]{article}
\usepackage[margin=25mm]{geometry}
\usepackage{fontspec}
\setmainfont{Latin Modern Roman}
\usepackage{amsmath,amssymb,mathtools,unicode-math}
\setmathfont{Latin Modern Math}
\usepackage{enumitem,needspace,xurl,hyperref}
\setlength{\parindent}{0pt}
\setlength{\parskip}{5pt}
\setlength{\emergencystretch}{3em}
\newcommand{\N}{\mathbb{N}}
\newcommand{\Z}{\mathbb{Z}}
\newcommand{\Q}{\mathbb{Q}}
\newcommand{\R}{\mathbb{R}}
\newcommand{\C}{\mathbb{C}}
\newcommand{\F}{\mathbb{F}}
\newcommand{\A}{\mathbb{A}}
\newcommand{\PP}{\mathbb{P}}
\newcommand{\E}{\mathbb{E}}
\newcommand{\eps}{\varepsilon}
\newcommand{\dd}{\,\mathrm{d}}
\begin{document}
\section*{UnsolvedMath ID: 30001397}
\textbf{Source problem code:} catalogue entry
\subsection{Definitions and mathematical statement}
An elliptic function is a nonconstant meromorphic function \(f\colon\C\to\widehat{\C}\) that is periodic with respect to a lattice of rank two. Write \(J(f)\) for its Julia set in the usual sense of meromorphic iteration. Assume that every critical point of \(f\) is attracted under iteration to an attracting or parabolic periodic cycle, and that \(f\) has at least one parabolic periodic point. An attracting cycle has multiplier of modulus less than one; a parabolic cycle has a root-of-unity multiplier, with a nonidentity return germ after a suitable iterate.

Let \(\mu\) denote the atomless \(h\)-conformal measure on \(J(f)\) referred to in the supplied source. Atomlessness means \(\mu(\{z\})=0\) for every \(z\in J(f)\). For a fixed conformal metric \(d\), the conformality convention is
\[
\mu(f(A))=\int_A \|Df(z)\|_d^{\,h}\,\mathrm d\mu(z)
\]
for Borel sets \(A\subseteq J(f)\) on which \(f\) is defined and injective. The exponent \(h\), the metric and derivative convention, and the normalization or selection of \(\mu\) are not specified by the supplied statement; this notation retains those unresolved choices rather than assigning them values.

Under a usual gauge convention, a gauge is a nondecreasing function \(g\colon[0,\infty)\to[0,\infty)\) with \(g(0)=0\), \(g(r)>0\) for \(r>0\), and \(\lim_{r\downarrow0}g(r)=0\). For a fixed metric \(d\), define
\[
\mathcal H_d^g(E)=\lim_{\delta\downarrow0}\inf\left\{\sum_{i=1}^{\infty}g\bigl(\operatorname{diam}_d U_i\bigr): E\subseteq\bigcup_{i=1}^{\infty}U_i,\ \operatorname{diam}_d U_i\leq\delta\right\}.
\]
The source does not specify its precise admissible gauge class or Hausdorff-measure normalization.

\textbf{Question, with these source-dependent choices retained.} For every elliptic function \(f\) satisfying the stated dynamical assumptions, does there exist a gauge \(g\), possibly depending on \(f\) and the designated measure \(\mu\), such that
\[
\mu(E)=\mathcal H_d^g(E)\qquad\text{for every Borel set }E\subseteq J(f)?
\]
Thus the requested conclusion is equality of measures on the Julia set. A fully specified formulation requires the missing conventions above, particularly the meaning of \(h\) and the intended measure \(\mu\).
\subsection{Short English statement}
For every elliptic map whose critical points approach attracting or parabolic cycles, with at least one parabolic cycle present, ask whether its designated atomless conformal measure on the Julia set is exactly a Hausdorff measure for some gauge. The supplied statement leaves the conformal exponent, metric, and measure normalization unspecified.
\subsection{Sources}
\begin{itemize}
\item Local UnsolvedMath registry, supplied complete record: permanent ID \(30001397\), source problem code \texttt{catalogue entry}, title \emph{Hausdorff Gauges for Conformal Measures of Elliptic Maps}; collection: Oberwolfach Reports Open Problems. No individual proposer is supplied; the proposed year is recorded as 2009.
\item Supplied report citation: \emph{Mini-Workshop: The Escaping Set in Transcendental Dynamics} (2010), Oberwolfach Reports, pp.\ 2927--2964, DOI \texttt{10.4171/owr/2009/54}. \url{https://doi.org/10.4171/owr/2009/54}. The background describes Oberwolfach Reports as documenting mathematical workshops and their open-problem discussions.
\item Supplied literature source, labelled ``Kotus Urbanski elliptic functions'': \url{https://doi.org/10.1007/BF02789304}. The supplied assessment describes this as a 2004 paper establishing the conformal measure and Hausdorff and packing measure behaviour, without constructing the requested gauge.
\item Supplied literature source, labelled ``original OWR report'': \url{https://doi.org/10.4171/owr/2009/54}.
\item Literature checked at: \texttt{2026-08-21}. This is supplied dated curation, not independently verified. The supplied assessment labels the problem open and says that the report's gauge-equivalence question remains unresolved. It explicitly describes this as dated literature triage derived from the supplied OWR curation and requests independent verification of time-sensitive status and cited claims.
\item The registry labels the statement \texttt{corrected\_verified} and reports that its mathematical content was checked against the official report, with notation, context, and question form normalized without changing the source-level claim. This verification description is supplied metadata, not verification performed for this draft.
\end{itemize}
\end{document}
